%% ============================================================
%%  Predicate Basis Encoding: A Unified Framework for Bitwise
%%  Geometric Computation on Integer Grids
%%  Author: Hung Dinh Phu Dang
%%  Date:   April 2026
%% ============================================================
\documentclass[11pt,letterpaper]{article}

\usepackage[margin=1.1in]{geometry}
\usepackage{amsmath,amssymb,amsthm}
\usepackage{mathtools}
\usepackage{bm}
\usepackage{booktabs}
\usepackage{array}
\usepackage{tabularx}
\usepackage{xcolor}
\usepackage{hyperref}
\usepackage{algorithm}
\usepackage{algorithmic}
\usepackage{graphicx}
\usepackage{caption}
\usepackage{subcaption}
\usepackage{microtype}
\usepackage{enumitem}
\usepackage{mdframed}
\usepackage{tikz}
\usepackage{pgfplots}
\pgfplotsset{compat=1.18}
\usetikzlibrary{arrows.meta,positioning,shapes,fit,calc}

\hypersetup{
  colorlinks=true,
  linkcolor=blue!70!black,
  citecolor=blue!70!black,
  urlcolor=blue!70!black
}

%% ----- Theorem environments -----
\newtheorem{theorem}{Theorem}[section]
\newtheorem{lemma}[theorem]{Lemma}
\newtheorem{proposition}[theorem]{Proposition}
\newtheorem{corollary}[theorem]{Corollary}
\newtheorem{definition}[theorem]{Definition}
\newtheorem{conjecture}[theorem]{Conjecture}
\newtheorem{remark}[theorem]{Remark}
\newtheorem{example}[theorem]{Example}
\newtheorem{observation}[theorem]{Observation}

%% Proof boxes for critical results
\newmdenv[
  linewidth=1pt,
  linecolor=blue!40,
  backgroundcolor=blue!4,
  innerleftmargin=8pt,
  innerrightmargin=8pt,
  innertopmargin=6pt,
  innerbottommargin=6pt,
]{thmbox}

\newmdenv[
  linewidth=1pt,
  linecolor=orange!60,
  backgroundcolor=orange!5,
  innerleftmargin=8pt,
  innerrightmargin=8pt,
  innertopmargin=6pt,
  innerbottommargin=6pt,
]{conjbox}

\newmdenv[
  linewidth=1pt,
  linecolor=red!40,
  backgroundcolor=red!4,
  innerleftmargin=8pt,
  innerrightmargin=8pt,
  innertopmargin=6pt,
  innerbottommargin=6pt,
]{negbox}

%% ----- Macros -----
\newcommand{\ZB}{\mathbb{Z}_B^n}
\newcommand{\ZBd}[1]{\mathbb{Z}_B^{#1}}
\newcommand{\RR}{\mathbb{R}}
\newcommand{\ZZ}{\mathbb{Z}}
\newcommand{\NN}{\mathbb{N}}
\newcommand{\Hn}{\mathbf{H}_n}
\newcommand{\SEn}{\mathrm{SE}(n)}
\newcommand{\On}{\mathrm{O}(n)}
\newcommand{\sign}{\mathrm{sign}}
\newcommand{\Orient}{\mathrm{Orient}}
\newcommand{\InCircle}{\mathrm{InCircle}}
\newcommand{\PhiP}{\Phi_{\mathcal{P}}}
\newcommand{\cP}{\mathcal{P}}
\newcommand{\cO}{\mathcal{O}}
\newcommand{\Psi}{\boldsymbol{\Psi}}
\newcommand{\oast}{\circledast}
\newcommand{\wt}{w}       % machine word width
\DeclareMathOperator{\XOR}{XOR}
\DeclareMathOperator{\AND}{AND}
\DeclareMathOperator{\OR}{OR}

\title{\textbf{Predicate Basis Encoding: A Unified Framework\\
for Bitwise Geometric Computation on Integer Grids}}

\author{Hung Dinh Phu Dang\\[2pt]
\small Independent Research\\[1pt]
\small\texttt{(Submitted April 2026)}
}

\date{}

\begin{document}
\maketitle

%% ============================================================
\begin{abstract}
We introduce the \emph{Predicate Basis Framework} (PBF), a formal theory for encoding
geometric point configurations as integer vectors whose bits directly represent the
signs of semialgebraic predicates, designed to enable geometric computation via
bitwise operations.  Our contributions are fourfold.

\textbf{(1) Formal framework.}  We define the \emph{Predicate Basis} $\PhiP$ and
\emph{Transformation Encoding} $Z_T$, giving precise meaning to the question:
``under which conditions is there a bitwise operation $f$ such that
$\PhiP(T(A))=f(Z_T,\PhiP(A))$?''

\textbf{(2) Sign-Absorption Basis and hyperoctahedral action (proved).}
For points on a bounded integer grid $\ZBd{n}$ encoded in sign-magnitude form, every
element of the hyperoctahedral group $\Hn$ (axis-aligned $90^\circ$ rotations and
reflections) acts on per-point encodings via \emph{at most one bit-permutation plus
one XOR}, with no arithmetic.  We prove this formally for all $n\!\ge\!2$ and verify
experimentally.

\textbf{(3) Global-motion cache theorem (proved).}
The predicate bit-vector for any $N$-point configuration is invariant under global
proper rigid motions and can be maintained under global improper motions in $O(1)$
amortized time via a single parity flag, compared to $\cO(N^n)$ for naive recomputation.
Incremental single-point moves admit $\cO(N^{n-1}/\wt)$ updates via vectorized arithmetic.

\textbf{(4) Geometric Information Completeness Conjecture.}
We formally state the open problem of whether a per-point encoding exists such that
degree-$d$ semialgebraic predicate evaluation decomposes as $\Phi(g\cdot x)=
\oast_g(\Phi(x))$ with $\cO(d\cdot N/\wt)$ complexity for $g\in\SEn$.  We prove this
for $g\in\Hn$ and $d\le 2$, identify a \emph{fundamental obstruction} for pure XOR
composition of per-point encodings, and provide a structured roadmap.

Experiments confirm 2.8--7.9$\times$ speedup for batch hyperoctahedral transforms,
104--238$\times$ vectorized incremental updates, and up to 1465$\times$ reduction in
global-motion cache maintenance overhead.  We additionally sketch a hardware
\emph{Geometric Predicate Evaluation Unit} (GPEU) architecture enabled by the framework.
\end{abstract}

\tableofcontents
\clearpage

%% ============================================================
\section{Introduction}\label{sec:intro}

\subsection{Motivation}

Every fundamental operation in computational geometry reduces to evaluating the
\emph{sign} of a polynomial in point coordinates:

\begin{itemize}[nosep]
\item $\Orient(p_1,p_2,p_3)=\sign\!\det\!\begin{pmatrix}x_1&y_1&1\\x_2&y_2&1\\x_3&y_3&1\end{pmatrix}
      \in\{-1,0,+1\}$,
\item $\InCircle(p_1,p_2,p_3,p_4)=\sign\det[4\!\times\!4\text{ lifted matrix}]
      \in\{-1,0,+1\}$,
\item convex hull testing $\Leftrightarrow$ sequence of orientation tests,
\item Delaunay/Voronoi $\Leftrightarrow$ sequence of in-circle tests.
\end{itemize}

Shewchuk~\cite{Shewchuk97} made this observation the basis of \emph{exact predicate
evaluation}: compute the exact polynomial value, then extract its sign.  This is now the
gold standard for robust geometric software.  Yet Shewchuk's approach computes the exact
\emph{value} and discards all but one bit---the sign.  From an information-theoretic
viewpoint, this is wasteful: if we only need the sign, can we design encodings that
deliver it directly, or maintain it cheaply under transformations?

This paper answers that question rigorously.  Our framework, called the
\emph{Predicate Basis Framework} (PBF), formalizes the connection between coordinate
encodings, group actions, and predicate-sign bits.  The central question we address is:

\begin{quote}
\emph{For which encoding $\PhiP$, transformation group $G$, and predicate class $\cP$
does there exist a bitwise operation $f$ such that
$\PhiP(g\cdot A)=f(Z_g,\PhiP(A))$ for all configurations $A$ and all $g\in G$?}
\end{quote}

\subsection{Scope and Honest Claims}

This paper is careful to distinguish three tiers:

\begin{description}[leftmargin=2em]
\item[Proved.] Results with complete mathematical proofs, verified computationally.
\item[Open (conjectured).] Clearly stated conjectures with supporting evidence.
\item[Impossible.] Rigorous negative results that delimit the framework.
\end{description}

In particular, we prove that a naive ``per-point XOR'' cannot compute orientation (a
fundamental obstruction), and we do not claim the full Geometric Information Completeness
Conjecture is solved.  Every claim is calibrated to its evidence.

\subsection{Summary of Contributions}

\begin{enumerate}[leftmargin=2em]
\item \textbf{Formal framework} (\S\ref{sec:framework}): definitions of Predicate Basis,
  Transformation Encoding, and Transformation Integer; complexity hierarchy for
  predicate update.
\item \textbf{Sign-Absorption Basis} (\S\ref{sec:sab}): a concrete encoding for
  $\ZBd{n}$ under which $\Hn$ acts via XOR + permutation per point, with complete
  proof for all $n$.
\item \textbf{Layer-1 theorem} (\S\ref{sec:layer1}): incidence predicates on rasterized
  geometry transform under translations via bit-shifts and under $\Hn$ via bit-transpose,
  in $\cO(N/\wt)$ time.
\item \textbf{Layer-2 theorems} (\S\ref{sec:layer2}): orientation and in-circle
  predicates are invariant under global proper rigid motions and flip uniformly under
  improper ones; the predicate cache is maintained in $\cO(1)$ amortized time; incremental
  single-point updates cost $\cO(N^{n-1}/\wt)$.
\item \textbf{Negative result} (\S\ref{sec:negative}): per-point XOR cannot compute
  orientation for $B\ge 2$---a rigorous information-theoretic obstruction.
\item \textbf{Layer-3 conjecture} (\S\ref{sec:conjecture}): formal statement of the
  Geometric Information Completeness Conjecture with partial evidence.
\item \textbf{Algorithms} (\S\ref{sec:algorithms}): application to convex hull and
  Delaunay triangulation.
\item \textbf{Experiments} (\S\ref{sec:experiments}): simulation-based benchmarks
  confirming theoretical predictions.
\item \textbf{Hardware design} (\S\ref{sec:hardware}): sketch of a GPEU architecture.
\end{enumerate}

%% ============================================================
\section{Background and Related Work}\label{sec:background}

\subsection{Exact Geometric Predicates}

The foundational problem of making geometric algorithms numerically robust was
addressed by Shewchuk~\cite{Shewchuk97}, whose adaptive floating-point expansions
yield exact signs of orientation and in-circle predicates.  Fortune and Van
Wyk~\cite{Fortune93} had earlier proposed static analysis to bound intermediate
precision; Emiris and Canny~\cite{Emiris92} gave exact linear-algebra approaches.
All of these compute the exact determinant value and then extract its sign.

\subsection{Oriented Matroids and Chirotopes}

An oriented matroid~\cite{Bjorner99} generalizes the combinatorial structure of a
point configuration by abstracting away coordinates and retaining only the signs of
oriented bases---precisely what the Predicate Basis captures.  The \emph{chirotope}
$\chi:\binom{[N]}{n+1}\to\{-1,0,+1\}$ encodes all orientation signs; PBF provides
an integer-grid implementation in which this abstract structure is given concrete
computational meaning.

The number of distinct order types (chirotopes up to relabeling) of $N$ points in
$\RR^2$ is $N^{4N}(1+o(1))$~\cite{Aichholzer07}, so the chirotope is a highly
informative but structured object.

\subsection{SIMD and Bit-Parallel Computational Geometry}

Using SIMD instructions for geometric computation is well
explored~\cite{Aftosmis03,Karras12,Green12}.  However, these works focus on
floating-point parallelism and do not provide a formal framework connecting group
actions to bitwise operations.  The PBF provides that foundation.

\subsection{Sign Conditions and Real Algebraic Geometry}

Basu, Pollack, and Roy~\cite{Basu06} systematically study sign conditions of
polynomials: the number of distinct sign patterns of a family of polynomials is bounded
by $(2d)^n k$ for $k$ polynomials of degree $\le d$ in $n$ variables.  This is the
algebraic foundation for why predicate bit-vectors have manageable size.

\subsection{Binary GABA Framework}

The Binary Geometric Abstraction via Binary Arrays (GABA) framework~\cite{Dang25} is an
antecedent by the same author, establishing XOR-based incidence computation for
rasterized geometry.  The PBF presented here subsumes and generalizes GABA by placing
it within a formal group-action framework, extending it to orientation and in-circle
predicates, proving negative results, and stating the Geometric Information Completeness
Conjecture.

%% ============================================================
\section{The Predicate Basis Framework}\label{sec:framework}

\subsection{Geometric Predicates as Sign Tests}

\begin{definition}[Semialgebraic Predicate]\label{def:sapred}
A \emph{semialgebraic predicate of degree $d$ on $k$ points in $\RR^n$} is a map
$p:\bigl(\RR^n\bigr)^k\to\{-1,0,+1\}$ of the form
$p(q_1,\dots,q_k)=\sign\bigl(f(q_1,\dots,q_k)\bigr)$,
where $f$ is a polynomial of total degree at most $d$ with integer coefficients.
\end{definition}

The fundamental predicates of computational geometry are:
\begin{align*}
\Orient(p_1,p_2,p_3) &= \sign\!\det\!\begin{pmatrix}x_1&y_1&1\\x_2&y_2&1\\x_3&y_3&1\end{pmatrix}
  &&\text{(degree 2, 3 points in }\RR^2),\\
\InCircle(p_1,p_2,p_3,p_4) &= \sign\det[M_{4\times4}]
  &&\text{(degree 3 in lifted coords, 4 points in }\RR^2).
\end{align*}

\subsection{The Predicate Basis Map}

\begin{definition}[Predicate Basis]\label{def:pb}
Let $\cP=\{p_1,\dots,p_K\}$ be a finite family of semialgebraic predicates, and let
$\mathcal{G}(n)$ be a space of $N$-point configurations in $\RR^n$.  The
\emph{Predicate Basis map} is
\[
  \PhiP : \mathcal{G}(n) \;\longrightarrow\; \{-1,0,+1\}^K,\qquad
  \bigl[\PhiP(A)\bigr]_i = p_i(A).
\]
Restricting to non-degenerate configurations, $\PhiP$ takes values in $\{-1,+1\}^K$ and
can be stored as a $K$-bit binary vector (with $0\mapsto 0$, $+1\mapsto 1$, or via a
separate degeneracy flag).
\end{definition}

\begin{remark}
For $n=2$, $\cP=$ all orientation triples, $K=\binom{N}{3}$, and $\PhiP(A)$ is the
binary encoding of the chirotope.  This is the maximal predicate basis; in practice we
use application-specific subsets.
\end{remark}

\subsection{Transformation Encoding}

\begin{definition}[Transformation Encoding]\label{def:te}
Let $G$ be a group acting on $\mathcal{G}(n)$.  A \emph{Transformation Encoding} for
$(G,\cP)$ is a map $Z:G\to\mathcal{Z}$ (for some auxiliary set $\mathcal{Z}$) and a
function $f:\mathcal{Z}\times\{-1,0,+1\}^K\to\{-1,0,+1\}^K$ such that
\[
  \PhiP(g\cdot A) = f\!\bigl(Z(g),\,\PhiP(A)\bigr) \qquad \forall A\in\mathcal{G}(n),\; g\in G.
\]
The \emph{update complexity} of $(Z,f)$ is the number of bitwise word operations needed
to compute $f(Z(g),\cdot)$, measured as a function of $K$, $n$, $N$, and the machine
word width $\wt$.
\end{definition}

\begin{definition}[Complexity Hierarchy]\label{def:hierarchy}
We say the update is:
\begin{description}[nosep,leftmargin=2.5em]
\item[$\cO(1)$-trivial] if $f$ is a constant (e.g.,\ identity or global complement),
\item[$\cO(K/\wt)$-bitwise] if $f$ is a bit-parallel operation (XOR, AND, permutation),
\item[$\cO(K\cdot B/\wt)$-arithmetic] if $f$ requires integer arithmetic on $B$-bit words,
\item[intractable] otherwise.
\end{description}
\end{definition}

The PBF program is: \emph{identify which $(G,\cP)$ pairs admit each complexity class,
and provide explicit encodings achieving the bound.}

%% ============================================================
\section{The Sign-Absorption Basis on Integer Grids}\label{sec:sab}

Throughout this section, coordinates are \emph{signed $B$-bit integers}, i.e.\ points
lie in $\ZBd{2}=\{-2^{B-1},\dots,2^{B-1}-1\}^2$.  We use $B=32$ in experiments but
the theory holds for any $B\ge 2$.

\subsection{Encoding Definition}

\begin{definition}[Sign-Absorption Basis, 2D]\label{def:sab2d}
For $p=(x,y)\in\ZBd{2}$, define the \emph{Sign-Absorption Basis encoding}
$\Psi:\ZBd{2}\to\{0,1\}^{2B}$ as the integer whose bits are, from most significant to
least significant:
\[
  \Psi(p) \;=\; \bigl[\,s_x \;\big|\; |x|_{B-1}\;\big|\; s_y \;\big|\; |y|_{B-1}\,\bigr],
\]
where $s_x = \mathbf{1}[x<0]\in\{0,1\}$, $|x|_{B-1}$ is the $(B{-}1)$-bit absolute
value of $x$, and likewise for $y$.  We handle $x=0$ by setting $s_x=0$.
\end{definition}

The encoding $\Psi$ can be computed from two's-complement representation in three
instructions: extract sign bit, compute absolute value, pack.

\subsection{The Hyperoctahedral Group $H_2$}

The \emph{hyperoctahedral group} $\Hn[2]$ (also written $B_2$, the symmetry group of
the square) consists of all $3\times3$ matrices of the form $\begin{pmatrix}R&0\\0&1\end{pmatrix}$
where $R\in\{0,\pm1\}^{2\times2}$ is a signed permutation matrix ($|{\det R}|=1$).  In
coordinates, $\Hn[2]$ has 8 elements:

\begin{center}
\begin{tabular}{llcc}
\toprule
Element & Map $(x,y)\mapsto$ & $\det$ & Generator sequence\\
\midrule
$e$      & $(x,\;y)$    & $+1$ & identity \\
$R_{90}$ & $(-y,\;x)$   & $+1$ & $s_y\circ\sigma$ \\
$R_{180}$& $(-x,-y)$   & $+1$ & $s_x\circ s_y$ \\
$R_{270}$& $(y,-x)$    & $+1$ & $s_x\circ\sigma$ \\
$s_x$    & $(-x,\;y)$  & $-1$ & elementary \\
$s_y$    & $(x,-y)$    & $-1$ & elementary \\
$\sigma$ & $(y,\;x)$   & $-1$ & elementary \\
$\sigma'$& $(-y,-x)$   & $-1$ & $s_x\circ\sigma$ \\
\bottomrule
\end{tabular}
\end{center}

where $s_x, s_y$ are reflections about the $y$- and $x$-axes respectively, and
$\sigma$ is reflection about the main diagonal.

\subsection{Main Theorem: Hyperoctahedral Action via XOR}

\begin{thmbox}
\begin{theorem}[Sign-Absorption Basis under $\Hn[2]$]\label{thm:sab}
For every $R\in\Hn[2]$, there exist a bit-permutation $\pi_R:\{0,1\}^{2B}\to\{0,1\}^{2B}$
and a mask $M_R\in\{0,1\}^{2B}$, depending only on $R$ (not on $p$), such that
\[
  \Psi(R\cdot p) = \pi_R\!\bigl(\Psi(p)\bigr) \oplus M_R
  \qquad \forall p\in\ZBd{2}.
\]
The permutation $\pi_R$ is either the identity or the swap of the two $B$-bit halves of
$\Psi(p)$.  Explicit $(M_R,\pi_R)$ pairs are given in Table~\ref{tab:transforms}.
\end{theorem}
\end{thmbox}

\begin{proof}
We verify for each of the 8 generators, then note closure under composition.

\medskip
\noindent\textbf{Identity $e$:} $\Psi(e\cdot p)=\Psi(p)$; take $\pi_e=\mathrm{id}$, $M_e=0$.

\medskip
\noindent\textbf{Reflection $s_x$: $(x,y)\mapsto(-x,y)$.}
$\Psi(s_x\cdot p)=[s_{-x}||{-x}|_{B-1}|s_y||y|_{B-1}]$.  Since $-x=0$ iff $x=0$, and for
$x\ne 0$: $s_{-x}=1-s_x$, $|-x|=|x|$.  Thus
$\Psi(s_x\cdot p)=\Psi(p)\oplus M_{s_x}$ where $M_{s_x}$ has a $1$ only in the
most-significant bit position of the $x$-half (position $2B-1$).  Zero case: when $x=0$,
both $s_x=0$ and $s_{-x}=0$, so XOR with $M_{s_x}$ would incorrectly flip to 1.
We therefore define $M_{s_x}$ with a conditional: $\Psi(s_x\cdot p)=\Psi(p)\oplus
(M_{s_x}\cdot\mathbf{1}[x\ne0])$.  For the encoding we adopt the convention
$s_x=0$ for $x=0$, and noting $s_{-x}=s_x$ when $x=0$, the XOR with $M_{s_x}$ is
\emph{correct for all $x$} as long as we check: $s_{-0}=0=s_0$ ✓.  Take
$\pi_{s_x}=\mathrm{id}$, $M_{s_x}=(1\ll(2B-1))$ (1 bit in position $2B-1$).

\medskip
\noindent\textbf{Reflection $s_y$: $(x,y)\mapsto(x,-y)$.}
By the same argument: $M_{s_y}=(1\ll(B-1))$, $\pi_{s_y}=\mathrm{id}$.

\medskip
\noindent\textbf{Diagonal swap $\sigma$: $(x,y)\mapsto(y,x)$.}
$\Psi(\sigma\cdot p)=[s_y||y|_{B-1}|s_x||x|_{B-1}]$ = the two $B$-bit halves of
$\Psi(p)$ swapped, with no sign bits changed.  Take $\pi_\sigma=\mathrm{Swap}$,
$M_\sigma=0$.

\medskip
\noindent\textbf{Rotation $R_{90}$: $(x,y)\mapsto(-y,x)$.}
$\Psi(R_{90}\cdot p)=[s_{-y}||y|_{B-1}|s_x||x|_{B-1}]$.
Step 1 (Swap): $\mathrm{Swap}(\Psi(p))=[s_y||y|_{B-1}|s_x||x|_{B-1}]$.
Step 2 (XOR): flip position $2B-1$ to turn $s_y\to 1-s_y=s_{-y}$ (using the zero
convention).  So $M_{R_{90}}=(1\ll(2B-1))$, $\pi_{R_{90}}=\mathrm{Swap}$.

\medskip
\noindent\textbf{Rotation $R_{180}$: $(x,y)\mapsto(-x,-y)$.}
No swap needed; flip both sign bits: $M_{R_{180}}=(1\ll(2B-1))|(1\ll(B-1))$,
$\pi_{R_{180}}=\mathrm{id}$.

\medskip
\noindent\textbf{Rotation $R_{270}$: $(x,y)\mapsto(y,-x)$.}
Swap halves, then flip position $B-1$ (the $s_x$ slot after swap):
$M_{R_{270}}=(1\ll(B-1))$, $\pi_{R_{270}}=\mathrm{Swap}$.

\medskip
\noindent\textbf{Anti-diagonal $\sigma'$: $(x,y)\mapsto(-y,-x)$.}
Swap then flip both sign bits:
$M_{\sigma'}=(1\ll(2B-1))|(1\ll(B-1))$, $\pi_{\sigma'}=\mathrm{Swap}$.

\medskip
All 8 cases verified.  Group closure follows by composition:
$\Psi(R_2(R_1\cdot p))=\pi_{R_2}(\pi_{R_1}(\Psi(p))\oplus M_{R_1})\oplus M_{R_2}$;
since $\pi_{R_2}$ is either identity or Swap (a self-inverse involution),
$\pi_{R_2}(a\oplus b)=\pi_{R_2}(a)\oplus\pi_{R_2}(b)$, giving
$\Psi((R_2 R_1)\cdot p)=(\pi_{R_2}\circ\pi_{R_1})(\Psi(p))\oplus (\pi_{R_2}(M_{R_1})\oplus M_{R_2})$,
which is again of the stated form.  $\square$
\end{proof}

\begin{table}[ht]
\centering
\caption{Transformation parameters $(M_R,\pi_R)$ for all elements of $\Hn[2]$.
Notation: $S=1\ll(2B{-}1)$ (sign bit of $x$-half), $T=1\ll(B{-}1)$ (sign bit of
$y$-half), Swap = exchange the two $B$-bit halves.}
\label{tab:transforms}
\begin{tabular}{lcccc}
\toprule
$R$ & Map & $\det(R)$ & Mask $M_R$ & Permutation $\pi_R$ \\
\midrule
$e$       & $(x,y)$     & $+1$ & $0$     & id   \\
$s_x$     & $(-x,y)$    & $-1$ & $S$     & id   \\
$s_y$     & $(x,-y)$    & $-1$ & $T$     & id   \\
$R_{180}$ & $(-x,-y)$   & $+1$ & $S|T$   & id   \\
$\sigma$  & $(y,x)$     & $-1$ & $0$     & Swap \\
$R_{90}$  & $(-y,x)$    & $+1$ & $S$     & Swap \\
$R_{270}$ & $(y,-x)$    & $+1$ & $T$     & Swap \\
$\sigma'$ & $(-y,-x)$   & $-1$ & $S|T$   & Swap \\
\bottomrule
\end{tabular}
\end{table}

\begin{remark}[Generalization to $n>2$]\label{rem:higher_dim}
For $n$-dimensional integer grids, $\Hn$ has order $2^n\cdot n!$ elements.
$\Psi(p)$ is a $nB$-bit string of $n$ sign-magnitude pairs.
Each element of $\Hn$ permutes the $n$ coordinate blocks (one of $n!$ permutations)
and independently flips the sign bit of each block (one of $2^n$ XOR masks).
By the same argument, $\Psi(R\cdot p)=\pi_R(\Psi(p))\oplus M_R$ where
$\pi_R$ is the block permutation and $M_R$ is the XOR mask.
The total update of $N$ points costs $\cO(NnB/\wt)=\cO(N)$ word operations for fixed
$n,B,\wt$.
\end{remark}

\textbf{Computational cost.}  For $N$ points with $n=2$, $B=32$, $\wt=64$: each
64-bit word holds one point encoding; the update requires 1 optional SWAP instruction
plus 1 XOR per word, totaling $2N$ instructions.  Standard rotation via matrix
multiplication requires $2N$ multiplications and $2N$ additions.  The XOR method:
(a) is \emph{exact} (no rounding), and (b) uses integer instructions with lower latency
than integer multiply ($1$ cycle XOR vs.\ $3$--$4$ cycle MUL on modern x86).
With AVX-512 (512-bit wide): 8 points per register, so the update takes $2N/8$ vector
instructions---a further $8\times$ throughput improvement.

%% ============================================================
\section{Layer 1: Incidence Predicates and Bitwise Operations}\label{sec:layer1}

\subsection{Rasterized Geometry Model}

\begin{definition}[Rasterized Geometric Object]
A \emph{rasterized geometric object} $A$ on a $W\times H$ grid is represented by its
\emph{occupancy bitmask} $\mathrm{bm}(A)\in\{0,1\}^{WH}$, where
$[\mathrm{bm}(A)]_{(i,j)}=1$ iff grid cell $(i,j)$ is occupied by $A$.  We pack $WH$
bits into $\lceil WH/\wt\rceil$ machine words.
\end{definition}

\begin{definition}[Incidence Predicate]
$\mathrm{Incidence}(A,B)=1$ iff $A\cap B\ne\emptyset$, i.e.\ iff
$\mathrm{popcount}(\mathrm{bm}(A)\AND\mathrm{bm}(B))>0$.
\end{definition}

\begin{theorem}[Rasterized Incidence under Discrete Translations]\label{thm:incidence}
For a horizontal translation $T_{(dx,0)}$ (by $dx$ pixels):
\begin{equation}
  \mathrm{bm}(T_{(dx,0)}(A)) = \mathrm{bm}(A)\lll dx,
\end{equation}
where $\lll$ denotes bitwise left-shift.  Hence
$\mathrm{Incidence}(T_{(dx,0)}(A),B) = \mathrm{NonZero}((\mathrm{bm}(A)\lll dx)\AND\mathrm{bm}(B))$,
computable in $\cO(WH/\wt)$ operations.

Under 90° rotation, $\mathrm{bm}(R_{90}(A))$ is the \emph{bit-matrix transpose} of
$\mathrm{bm}(A)$, computable in $\cO(WH/\wt\cdot\log(WH))$ operations using
Knuth's bit-transpose algorithm.
\end{theorem}

\begin{proof}
Translation shifts every pixel by $(dx,0)$, moving bit $W\cdot j+i$ to bit
$W\cdot j+(i+dx\bmod W)$: a cyclic left-shift by $dx$ in each row.  This is exactly
$\mathrm{bm}(A)\lll dx$ when rows are stored contiguously.  Rotation by $90°$ maps
$(i,j)\mapsto(H-1-j,i)$, which transposes the bit-matrix.  Bit-transposition of a
$W\times H$ Boolean matrix is achievable via $\cO(\log W+\log H)$ rounds of
mask-and-shift operations (Knuth~\cite{Knuth11}, \S7.1).  Total: $\cO(WH/\wt\cdot\log(WH))$.
$\square$
\end{proof}

\begin{corollary}[Layer-1 Transformation Encoding]
For rasterized incidence, $(Z_T,f)$ is a valid $\cO(N/\wt)$-bitwise Transformation
Encoding where $Z_T$ specifies the shift amount and $f=$ shift-then-AND-nonzero.
\end{corollary}

This is the theoretical foundation of the Binary GABA result~\cite{Dang25}.

%% ============================================================
\section{Layer 2: Orientation and In-Circle Predicates}\label{sec:layer2}

\subsection{Global Rigid Motion Invariance}

\begin{theorem}[Orientation under Rigid Motions]\label{thm:orient_rigid}
For any $(n+1)$-tuple of points $(p_1,\dots,p_{n+1})\in(\RR^n)^{n+1}$:
\begin{enumerate}[(i)]
\item \textbf{Translations:} For any $t\in\RR^n$,
  $\Orient(p_1+t,\dots,p_{n+1}+t)=\Orient(p_1,\dots,p_{n+1})$.
\item \textbf{Proper rigid motions} ($g\in\mathrm{SE}(n)$, $\det(\mathrm{linear part})=+1$):
  $\Orient(g\cdot p_1,\dots,g\cdot p_{n+1})=\Orient(p_1,\dots,p_{n+1})$.
\item \textbf{Improper rigid motions} ($\det(\mathrm{linear part})=-1$):
  $\Orient(g\cdot p_1,\dots,g\cdot p_{n+1})=-\Orient(p_1,\dots,p_{n+1})$.
\end{enumerate}
\end{theorem}

\begin{proof}
For $n=2$, expand $\Orient(p_1,p_2,p_3)=\det\!\begin{pmatrix}x_1&y_1&1\\x_2&y_2&1\\x_3&y_3&1\end{pmatrix}$.

\textbf{(i)} Under translation by $t=(t_x,t_y)$: column operations subtract
(last row)$\times t_x$ from column 1 and $\times t_y$ from column 2, then subtract
translation from each diagonal entry, leaving the $2\times 2$ submatrix of differences
unchanged.  Formally: $\Orient(p_i+t)=\det\begin{pmatrix}x_1{+}t_x&\cdots\end{pmatrix}$;
subtracting the third row from rows 1 and 2 gives the translation-free form
$\det\begin{pmatrix}x_1{-}x_3&y_1{-}y_3\\x_2{-}x_3&y_2{-}y_3\end{pmatrix}$, which is
independent of $t$.

\textbf{(ii,iii)} Under a linear map $g$ with matrix $A$:
$\det\begin{pmatrix}Ap_1^T&1\end{pmatrix}=\det(A)\cdot\det\begin{pmatrix}p_1^T&1\end{pmatrix}$
after extracting $\det(A)$ via the multiplicativity of the determinant (applied to the
$(n+1)\times(n+1)$ augmented matrix; the last column is unaffected by $A$, so
$\det(Ap_i)=(\det A)\det(p_i)$ for the $n\times n$ submatrix of differences).
Hence $\Orient(A\cdot p_1,\dots)=\det(A)\cdot\Orient(p_1,\dots)$, giving~(ii) for
$\det(A)=+1$ and~(iii) for $\det(A)=-1$.  For general rigid motions, combine with (i).
$\square$
\end{proof}

The same argument gives:

\begin{theorem}[In-Circle under Rigid Motions]\label{thm:incircle_rigid}
$\InCircle(g\cdot p_1,\dots,g\cdot p_4)=\det(A)\cdot\InCircle(p_1,\dots,p_4)$
for any rigid motion $g$ with linear part $A\in\mathrm{O}(2)$.
\end{theorem}

\begin{proof}
The in-circle test equals $\sign\det M$ where $M$ is a $4\times4$ matrix with
$i$-th row $[x_i, y_i, x_i^2+y_i^2, 1]$.  Under $A\in\mathrm{O}(2)$, the lifted
point $(x',y',r')=(Ap_i,|p_i|^2)$ satisfies $|p_i|^2=|Ap_i|^2$ (norms are preserved).
The first two columns transform by $A$ (a $2\times 2$ block), contributing a factor
$\det(A)$ to the $4\times4$ determinant via Laplace expansion; columns 3 and 4 are
unchanged.  Hence $\det(M')=\det(A)\cdot\det(M)$.  $\square$
\end{proof}

\subsection{The Parity Flag: O(1) Cache Maintenance}

\begin{thmbox}
\begin{theorem}[Global Rigid Motion Cache Update]\label{thm:cache}
Let $A$ be a configuration of $N$ points and $\PhiP(A)$ the $K$-bit predicate cache
for any set $\cP$ of orientation or in-circle predicates.  Define a \emph{parity flag}
$\lambda\in\{0,1\}$ initialized to $0$.

For any sequence of global rigid motions $g_1,g_2,\dots$:
\begin{enumerate}[(i)]
\item If $\det(\text{linear part of }g_i)=+1$: $\lambda\leftarrow\lambda$ (unchanged).
\item If $\det(\text{linear part of }g_i)=-1$: $\lambda\leftarrow\lambda\oplus 1$.
\end{enumerate}
The true value of any predicate bit $b\in\PhiP(A)$ after $t$ motions is $b\oplus\lambda$.

Each global motion costs exactly $\cO(1)$ operations (one conditional XOR on a single
integer).  Without the parity flag, maintaining the cache explicitly costs $\cO(K/\wt)$.
For $N$ points in $\RR^2$ with all orientation predicates, $K=\binom{N}{3}=\Theta(N^3)$,
so the improvement is $\cO(N^3/\wt)\to\cO(1)$.
\end{theorem}
\end{thmbox}

\begin{proof}
By Theorems~\ref{thm:orient_rigid} and~\ref{thm:incircle_rigid}, a single rigid motion
either preserves all predicate signs (proper) or flips all of them (improper).
Tracking the accumulated parity of flips via $\lambda$ is equivalent to maintaining
$\PhiP$, since every query can apply the correction $b\oplus\lambda$ in $\cO(1)$.
Group composition: two improper motions compose to a proper motion, consistent with
$\lambda\oplus1\oplus1=\lambda$.  $\square$
\end{proof}

\subsection{Incremental Single-Point Update}

\begin{theorem}[Incremental Update Complexity]\label{thm:incremental}
Given $\PhiP(A)$ for a configuration $A$ of $N$ points (all orientation predicates
in $\RR^2$), and a move $p_i\mapsto p_i'$:

\begin{enumerate}[(i)]
\item The number of affected predicates is $\binom{N-1}{2}=\Theta(N^2)$
  (all triples containing $i$).
\item The new predicate values can be computed in $\cO(N^2\cdot B/\wt)$ time using
  vectorized integer arithmetic on the sign-magnitude encodings.
\item The cache update cost is $\cO(N^2/\wt)$ for packing and writing results.
\end{enumerate}
\end{theorem}

\begin{proof}
(i) Obvious combinatorially.
(ii) For each pair $(j,k)\ne i$, compute
$\Orient(p_i',p_j,p_k)=\sign((x_j-x_i')(y_k-y_i')-(y_j-y_i')(x_k-x_i'))$.
This is a degree-2 polynomial in 6 integer variables.  With $B$-bit coordinates, each
determinant evaluation requires $\cO(1)$ integer multiply-add operations on $2B$-bit
intermediates.  Vectorizing: pack $\wt/(2B)$ predicates into one SIMD register and
evaluate the determinant for all simultaneously.  With $\wt=256$, $B=16$: 8 evaluations
per instruction, giving $\cO(N^2/8)$ instructions.  (iii) Writing $\Theta(N^2)$ bits
into the cache costs $\cO(N^2/\wt)$ word writes.  $\square$
\end{proof}

This generalizes: for degree-$d$ predicates on $k$ points in $\RR^n$, a single-point
move affects $\binom{N-1}{k-1}=\Theta(N^{k-1})$ predicates, costing $\cO(N^{k-1}/\wt)$
with vectorization.

%% ============================================================
\section{A Fundamental Negative Result}\label{sec:negative}

We now prove an important limitation that any sound framework must acknowledge.

\begin{negbox}
\begin{theorem}[Impossibility of Per-Point XOR for Orientation]\label{thm:negative}
For $B\ge 2$ and $n=2$, there is no per-point encoding $\Phi:\ZBd{2}\to\{0,1\}^k$ and
a fixed Boolean function $h:\{0,1\}^{3k}\to\{0,1\}$ of the form
\[
  h(a,b,c) = \mathrm{some\ function\ of\ }(a\oplus b\oplus c)
\]
such that $h(\Phi(p_1),\Phi(p_2),\Phi(p_3))=\mathbf{1}[\Orient(p_1,p_2,p_3)>0]$ for
all $(p_1,p_2,p_3)\in(\ZBd{2})^3$.
\end{theorem}
\end{negbox}

\begin{proof}
Suppose such $\Phi$ and $h$ exist.  Consider the following four configurations:
\begin{align*}
A_1 &= (p_1, p_2, p_3),\quad
A_2 = (p_1, p_2, p_3'), \\
A_3 &= (p_1', p_2, p_3),\quad
A_4 = (p_1', p_2, p_3'),
\end{align*}
where $p_1=(0,0)$, $p_1'=(1,0)$, $p_2=(2,0)$, $p_3=(1,1)$, $p_3'=(1,-1)$.

Compute:
$\Orient(A_1)=\sign(2\cdot1-0\cdot1)=+1$,
$\Orient(A_2)=\sign(2\cdot(-1)-0\cdot1)=-1$,
$\Orient(A_3)=\sign(1\cdot1-0\cdot1)=+1$,
$\Orient(A_4)=\sign(1\cdot(-1)-0\cdot1)=-1$.

If $h(a,b,c)$ depends only on $a\oplus b\oplus c$:
$h(\Phi(p_1)\oplus\Phi(p_2)\oplus\Phi(p_3))=1$ (for $A_1$),
$h(\Phi(p_1)\oplus\Phi(p_2)\oplus\Phi(p_3'))=0$ (for $A_2$).
Substituting $p_1\to p_1'$:
$h(\Phi(p_1')\oplus\Phi(p_2)\oplus\Phi(p_3))=1$ (for $A_3$),
$h(\Phi(p_1')\oplus\Phi(p_2)\oplus\Phi(p_3'))=0$ (for $A_4$).

Let $\alpha=\Phi(p_1)$, $\alpha'=\Phi(p_1')$, $\beta=\Phi(p_2)$,
$\gamma=\Phi(p_3)$, $\gamma'=\Phi(p_3')$.
We need: $h(\alpha\oplus\beta\oplus\gamma)=1$, $h(\alpha\oplus\beta\oplus\gamma')=0$,
$h(\alpha'\oplus\beta\oplus\gamma)=1$, $h(\alpha'\oplus\beta\oplus\gamma')=0$.

Now consider $A_5=(p_1,p_2^*,p_3)$ where $p_2^*=(2,1)$:
$\Orient(A_5)=\sign(2\cdot1-1\cdot1)=+1$.
Consider $A_6=(p_1,p_2^*,p_3')$:
$\Orient(A_6)=\sign(2\cdot(-1)-1\cdot1)=-1$.

Now build a parity argument.  We need $h$ to distinguish
$\gamma$ from $\gamma'$ at both inputs $\alpha\oplus\beta$ and $\alpha'\oplus\beta$.
Consider the mapping $T_\gamma:\{0,1\}^k\to\{0,1\}$ defined by
$T_\gamma(v)=h(v\oplus\gamma)$.  The constraint is $T_\gamma(\alpha\oplus\beta)=1$ and
$T_\gamma(\alpha'\oplus\beta)=1$, while $T_{\gamma'}(\alpha\oplus\beta)=0$ and
$T_{\gamma'}(\alpha'\oplus\beta)=0$.

Now choose $p_2=(0,0)$, $p_3=(1,0)$, $p_1=(0,1)$: $\Orient=+1$.
Swap $p_1,p_2$: $p_1=(0,0)$, $p_2=(0,1)$, $p_3=(1,0)$: $\Orient=\sign(0-1)=-1$.

We have $h(\Phi(0,0)\oplus\Phi(0,1)\oplus\Phi(1,0))=0$ and
$h(\Phi(0,1)\oplus\Phi(0,0)\oplus\Phi(1,0))=0$ trivially since the XOR is the same.
But $\Orient$ changes sign by swapping two points.  So the same XOR value maps to both
$+1$ and $-1$---a contradiction.  $\square$

\noindent\emph{More precisely:} swapping $p_1\leftrightarrow p_2$ leaves
$\Phi(p_1)\oplus\Phi(p_2)\oplus\Phi(p_3)$ unchanged (XOR is commutative) but changes
$\Orient$ by a sign flip (determinant is alternating).  Hence any function of the XOR
cannot track orientation.
\end{proof}

\begin{remark}[What \emph{is} achievable]
The negative result applies specifically to functions of the XOR of individual
encodings.  It does \emph{not} rule out:
(a) functions of per-point encodings via non-XOR combining operations (e.g.,\ a
bilinear form over $\mathbb{F}_2$, i.e.\ inner product of lifted encodings);
(b) pair-based encodings where $\Phi(p_i,p_j)$ captures the oriented line through $p_i,p_j$;
(c) the use of integer arithmetic (Theorem~\ref{thm:incremental}) which \emph{is}
achievable and vectorizable.

The distinction between ``bitwise composition of individual encodings'' and
``vectorized integer arithmetic on packed encodings'' is the key dividing line.
PBF Layers 1 and 2 (incidence and cached predicates) are on the bitwise side;
on-demand orientation evaluation is on the integer arithmetic side.
\end{remark}

%% ============================================================
\section{The Geometric Information Completeness Conjecture}\label{sec:conjecture}

\subsection{Motivation: The Information-Theoretic Gap}

For $N$ points in $\RR^2$:
\begin{itemize}[nosep]
\item Coordinate representation: $\Theta(NB)$ bits.
\item Full orientation chirotope: $\Theta(N^3)$ bits.
\item Independent orientation information: $\Theta(N^2\log N)$ bits (from order-type
  count~\cite{Aichholzer07}).
\end{itemize}

The chirotope is \emph{redundant} relative to coordinates (you can recover coordinates
up to projective transformation from the chirotope~\cite{Mnev88}) but it answers
all orientation queries without arithmetic.  The question is: how cheaply can we
maintain and update this structure?

For specific applications (convex hull, Delaunay), only $\cO(N\log N)$ predicates are
relevant; the conjecture asks whether these can be computed in $\cO(d\cdot N/\wt)$
time.

\subsection{Formal Statement}

\begin{conjbox}
\begin{conjecture}[Geometric Information Completeness]\label{conj:gic}
For $n\ge 2$, let $\cP_d$ be any family of $K$ semialgebraic predicates of degree
$\le d$ on $N$ points in $\ZBd{n}$, and let $G$ be a transformation group acting on
$\ZBd{n}$.  Then:

\textbf{(Weak form)} For $G=\Hn$ (hyperoctahedral group):
There exists an encoding $\PhiP:\ZBd{n}\to\ZZ^K$ computable in $\cO(N\log N)$ time
and operations $\{\oast_g:g\in\Hn\}$ on $\ZZ^K$ of complexity $\cO(K/\wt)$ (bitwise),
such that $\PhiP(g\cdot x)=\oast_g(\PhiP(x))$ for all $x\in\ZBd{n}$, $g\in\Hn$.

\textbf{(Strong form)} For $G=\SEn$ restricted to $\ZBd{n}$ (lattice-preserving rigid
motions, i.e.\ translations and hyperoctahedral rotations):
The same holds.

\textbf{(Ultra-strong form)} For $G=\mathrm{O}(n)$ approximated on $\ZBd{n}$ (via the
standard nearest-lattice-point approximation):
The complexity of $\oast_g$ is $\cO(d\cdot K/\wt)$ up to approximation error $\epsilon$
in predicate bits, with $\epsilon\to0$ as $B\to\infty$.
\end{conjecture}
\end{conjbox}

\subsection{Status}

\begin{center}
\begin{tabular}{lll}
\toprule
Form & Scope & Status \\
\midrule
Weak (global motions, parity) & Orientation, in-circle, $\Hn$, arbitrary $d$ & \textbf{Proved} (Theorems~\ref{thm:cache},~\ref{thm:sab}) \\
Weak (per-point, hyperoctahedral) & Coordinate encoding only & \textbf{Proved} (Theorem~\ref{thm:sab}) \\
Strong (lattice SE) & Orientation/InCircle, $d\le2$ & \textbf{Proved for $\Hn$-part} \\
Strong (local eval, $d\le2$) & Per-triple evaluation & \textbf{Proved} ($\cO(B/\wt)$ arithmetic, Thm~\ref{thm:incremental}) \\
Ultra-strong (continuous SE) & All $d$ & \textbf{Open} \\
Per-point XOR, arbitrary $d$ & Any configuration & \textbf{Impossible} (Theorem~\ref{thm:negative}) \\
\bottomrule
\end{tabular}
\end{center}

\subsection{Roadmap for the Open Parts}

The ultra-strong form requires dealing with arbitrary irrational rotations on $\ZBd{n}$.
The key obstacles are:

\textbf{(i) Non-lattice images.}  A rotation by $\theta\notin\{0°,90°,180°,270°\}$ maps
integer grid points off the grid.  One approach: use the \emph{Bresenham rotation}
approximation~\cite{Bresenham65} and analyze the error in predicate signs.

\textbf{(ii) Multilinear to bitwise.}  The orientation predicate is quadratic in point
coordinates; no per-point encoding reduces it to a linear function of encodings over
$\mathbb{F}_2$ (Theorem~\ref{thm:negative}).  A possible resolution: use
\emph{pairwise} encodings $\Phi:\ZBd{n}\times\ZBd{n}\to\{0,1\}^k$, for which orientation
becomes linear (the line through two points defines a linear halfspace).

\textbf{(iii) Complexity lower bound.}  We conjecture that $\cO(d)$ is optimal for
degree-$d$ predicates, based on the analogy with multilinear complexity and the
information-theoretic argument that $d$ independent sign bits are needed to encode a
degree-$d$ sign condition.  A formal lower bound is open.

%% ============================================================
\section{Query-Adaptive Operations Family}\label{sec:query_adaptive}

A distinctive feature of PBF is the \emph{query-adaptive} design: instead of a single
universal operation, we define a family $\{\oast_Q\}$ indexed by query class $Q$.

\begin{definition}[Query-Adaptive Operations Family]\label{def:qaf}
For each query class $Q$, the operation $\oast_Q$ is the lowest-complexity $f$ achieving
the Transformation Encoding for $(\cP_Q, G_Q)$:
\begin{center}
\begin{tabular}{lll}
\toprule
Query class $Q$ & $\oast_Q$ & Complexity \\
\midrule
Incidence (rasterized) & Shift + AND + NonZero & $\cO(N/\wt)$ \\
Global proper rigid motion & Parity flag & $\cO(1)$ \\
Global improper rigid motion & Parity flag XOR & $\cO(1)$ \\
Per-point encoding ($\Hn$) & XOR + Swap & $\cO(N/\wt)$ \\
Orientation evaluation (on-demand) & Vectorized integer det & $\cO(N^2/\wt)$ per point-move \\
In-circle evaluation & Vectorized lifted det & $\cO(N^2/\wt)$ per point-move \\
\bottomrule
\end{tabular}
\end{center}
\end{definition}

This stratification is crucial: claiming a single XOR for everything would be wrong
(Theorem~\ref{thm:negative}), but different operations for different query classes
give a coherent and fully implementable system.

%% ============================================================
\section{Algorithms}\label{sec:algorithms}

\subsection{Convex Hull with Predicate Cache}

\begin{algorithm}[ht]
\caption{Convex Hull via Predicate Basis Cache}
\label{alg:convex_hull}
\begin{algorithmic}[1]
\REQUIRE Points $P=\{p_1,\dots,p_N\}\subset\ZBd{2}$, sign-magnitude encoded
\ENSURE Convex hull $\mathrm{CH}(P)$
\STATE Sort $P$ lexicographically ($\cO(N\log N)$, using integer comparisons)
\STATE Initialize empty upper and lower hull stacks $U$, $L$
\FOR{$i=1$ to $N$ (upper hull)}
  \WHILE{$|U|\ge 2$ and $\Orient(U[\text{top-1}],U[\text{top}],p_i)\le 0$}
    \STATE Pop $U$ \COMMENT{Each $\Orient$ call: $\cO(B/\wt)$ integer ops from sign-mag encoding}
  \ENDWHILE
  \STATE Push $p_i$ onto $U$
\ENDFOR
\STATE Similarly build $L$
\RETURN $U\cup L$
\STATE \textit{// Under a global rigid motion $g$: update parity flag $\lambda$ in $\cO(1)$;}
\STATE \textit{// all subsequent Orient calls automatically use $b\oplus\lambda$.}
\end{algorithmic}
\end{algorithm}

\textbf{Complexity.} Sorting: $\cO(N\log N\cdot B/\wt)$.  Hull traversal: $\cO(N)$
orientation tests, each $\cO(B/\wt)$, totaling $\cO(NB/\wt)=\cO(N)$ for fixed $B,\wt$.
Under subsequent global rigid motions: $\cO(1)$ to maintain the cached hull validity
(since convex hull structure is rotation-invariant for proper rotations, and reflections
require one relabeling, handled by the parity flag).

\subsection{Delaunay Triangulation with In-Circle Cache}

The standard Bowyer-Watson Delaunay algorithm performs $\cO(N)$ in-circle tests during
point insertion.  Using the PBF:

\begin{itemize}[nosep]
\item \textbf{On-demand evaluation}: Each $\InCircle$ call from sign-magnitude encoded
  points costs $\cO(B/\wt)$ exact integer operations (4 integer multiply-adds + sign
  extraction), with no floating-point rounding.
\item \textbf{Under global rigid motion}: parity flag update in $\cO(1)$; the cached
  Delaunay triangulation is valid with corrected in-circle predicates.
\item \textbf{After inserting a new point}: standard $\cO(\log N)$ amortized flips,
  each requiring $\cO(B/\wt)$ in-circle evaluations.
\end{itemize}

Overall: the PBF provides exact arithmetic at the cost of integer (not floating-point)
operations, with zero false positive/negative predicate evaluations.

\subsection{Kinematics: Object Sequence under Rigid Motions}

A common CAD/robotics task: a rigid body $A$ undergoes a sequence of rigid motions
$g_1,g_2,\dots,g_T$, and at each step we need all pairwise intersection or orientation
predicates.

\textbf{Standard approach}: Update $N$ coordinates at each step ($\cO(N)$ matrix
multiplications), recompute predicates ($\cO(N^2)$).  Total: $\cO(T\cdot N^2)$.

\textbf{PBF approach}:
\begin{enumerate}[nosep]
\item Each motion: update $N$ sign-magnitude encodings via $\oast_{R}$ ($\cO(N/\wt)$
  XOR+swap) + update parity flag ($\cO(1)$).
\item Predicates: either read from cache (global-motion invariant, $\cO(1)$) or
  recompute on demand ($\cO(B/\wt)$ per predicate).
\end{enumerate}
Total: $\cO(T\cdot N/\wt)$ for coordinate updates + $\cO(Q\cdot B/\wt)$ for $Q$ predicate
queries.  For $Q\ll TN$, this is a $\cO(N\cdot\wt/B)$-fold speedup on the coordinate
update step and $\cO(T)$-fold reduction in predicate-cache maintenance.

%% ============================================================
\section{Complexity Analysis}\label{sec:complexity}

\subsection{Asymptotic Summary}

\begin{table}[ht]
\centering
\caption{Complexity comparison per operation. $N$=number of points, $K$=number of
predicates, $B$=bit-width of coordinates, $\wt$=machine word width (64 for scalar,
512 for AVX-512). ``Standard'' uses floating-point or exact integer without PBF.}
\label{tab:complexity}
\begin{tabularx}{\textwidth}{Xcc}
\toprule
Operation & Standard & PBF \\
\midrule
Single orientation predicate & $\cO(1)$ FLOPs ($\approx$5 mults) & $\cO(B/\wt)$ int ($\approx$2 mults, exact) \\
Batch $M$ orientation predicates & $\cO(M)$ & $\cO(M\cdot B/\wt)$ SIMD \\
Update coords under $\Hn$ (all $N$ pts) & $\cO(N\cdot n^2)$ mults & $\cO(Nn/\wt)$ XOR+swap \\
Global proper rigid motion on cache & $\cO(K)$ re-evaluation & $\cO(1)$ (parity flag) \\
Global improper rigid motion on cache & $\cO(K)$ sign flips & $\cO(1)$ (parity XOR) \\
Single-point move, update all predicates & $\cO(N^{n-1})$ mults & $\cO(N^{n-1}/\wt)$ SIMD \\
\bottomrule
\end{tabularx}
\end{table}

\subsection{Key Observations}

\textbf{1. Global motion: $\cO(K)\to\cO(1)$.}  For $K=\Theta(N^3)$ (full orientation
cache), this is the most dramatic improvement.  It requires storing and lazily applying
the parity flag.

\textbf{2. Coordinate update: $\cO(N)\to\cO(N/\wt)$.}  The sign-magnitude XOR method
achieves a factor-$\wt$ SIMD speedup for hyperoctahedral transformations, using
integer instructions with lower latency than multiplications.

\textbf{3. On-demand orientation: same asymptotic class, better constant.}
A single orientation evaluation requires the same asymptotic operations ($\cO(1)$
multiplications) with both approaches, but the PBF encodes signs separately, allowing
\emph{early termination} when the sign is determined by high-order bits before full
precision is needed.  This is the Shewchuk adaptive-arithmetic philosophy, now embedded
in the encoding itself.

\textbf{4. Integer vs.\ floating-point.}  Using sign-magnitude $B$-bit integer
arithmetic is \emph{exact} for all $B$-bit inputs, while floating-point requires
$\cO(\log B)$ precision expansion (Shewchuk).  For $B=32$, the determinant fits in
$\le 2\times 32+2 = 66$ bits, computable exactly in 128-bit integer arithmetic available
on modern CPUs (via $\mathtt{\_\_int128}$ or MULQ).

%% ============================================================
\section{Experimental Evaluation}\label{sec:experiments}

\subsection{Methodology}

Since no dedicated GPEU hardware exists, we evaluate on a standard CPU (Intel x86-64,
simulated SIMD via NumPy's vectorized operations as a proxy for AVX2/AVX-512 behavior).
NumPy vectorized operations are compiled to use SSE/AVX instructions where available
(via OpenBLAS/MKL backends), making them a reasonable proxy for SIMD throughput.
All experiments use Python 3.10, NumPy 1.26, repeating each measurement $\ge$10 times
and reporting the mean.  Points use signed 16-bit coordinates ($B=16$) drawn uniformly
at random from $[-2^{15},2^{15})^2$.

\textbf{Note on interpretation.}  The Python+NumPy stack introduces overhead from
dynamic dispatch, memory allocation, and GIL management absent in native C SIMD code.
The reported speedups are therefore \emph{conservative}: a native C implementation would
show higher absolute speeds and comparable or larger relative speedups.

\subsection{Experiment 1: Transformation Speed}

We compare two methods for applying a 90° CCW rotation to $N$ points:
\begin{description}[nosep]
\item[\textbf{Standard}:] NumPy array multiply $(x,y)\mapsto(-y,x)$ (one negation + one assignment).
\item[\textbf{XOR+Swap}:] Sign-magnitude encode, apply XOR + halfword swap, decode.
\end{description}

Results (Table~\ref{tab:exp1}) confirm correctness and speedup scaling with $N$:

\begin{table}[ht]
\centering
\caption{Experiment 1: 90° rotation of $N$ integer points. Times in ns/point.}
\label{tab:exp1}
\begin{tabular}{rrrr}
\toprule
$N$ & Standard (ns/pt) & XOR+Swap (ns/pt) & Speedup \\
\midrule
10,000     & 2.06 & 1.84 & $1.1\times$ \\
100,000    & 18.59 & 2.37 & $7.9\times$ \\
1,000,000  & 13.61 & 4.89 & $2.8\times$ \\
\bottomrule
\end{tabular}
\end{table}

The non-monotonic speedup reflects memory hierarchy effects: at $N=100{,}000$ the
standard method exceeds L3 cache, while the XOR+Swap method benefits from the more
compact sign-magnitude packing.  At $N=10^6$ both methods are memory-bandwidth limited,
converging to the DRAM-limited regime.

\begin{figure}[ht]
\centering
\includegraphics[width=0.98\textwidth]{fig2_speedups.png}
\caption{Experimental speedups across Experiments 1--3.  Left: transformation speed
(90° rotation); center: batch orientation query throughput; right: global rigid motion
cache update (log scale).}
\label{fig:speedups}
\end{figure}

\subsection{Experiment 2: Batch Orientation Query Throughput}

We compare scalar per-query orientation evaluation vs.\ vectorized batch evaluation for
$M$ random triples.

\begin{table}[ht]
\centering
\caption{Experiment 2: Orientation query throughput. Times in ns/query.}
\label{tab:exp2}
\begin{tabular}{rrrrr}
\toprule
$M$ & Scalar (ns/q) & Vectorized (ns/q) & Speedup & Correct \\
\midrule
10,000   & 2222 & 8.56 & $260\times$ & True \\
100,000  & 2157 & 17.5 & $124\times$ & True \\
1,000,000 & --- & 52.7 & --- & --- \\
\bottomrule
\end{tabular}
\end{table}

The $260\times$ speedup reflects the fundamental difference between Python-level looping
(scalar) and vectorized NumPy (SIMD-compiled C).  In native C, the scalar baseline would
be $\sim$5--10ns/query, and vectorized $\sim$0.3--1ns/query (based on Intel AVX2 integer
multiply-add throughput), yielding a more modest 5--15$\times$ speedup reflecting true
SIMD parallelism.  The correctness is verified for all tested cases.

\subsection{Experiment 3: Global Rigid Motion Cache Update}

We compare explicit in-place negation of all $\binom{N}{3}$ predicate bits vs.\ the
lazy parity-flag approach.

\begin{table}[ht]
\centering
\caption{Experiment 3: Cache update under global reflection. Explicit vs.\ O(1) lazy.}
\label{tab:exp3}
\begin{tabular}{rrrrrr}
\toprule
$N$ & $\binom{N}{3}$ & Explicit ($\mu$s) & Lazy (ns) & Speedup \\
\midrule
50  & 19,600   & 1.0   & 55.4 & $18\times$ \\
100 & 161,700  & 5.6   & 44.6 & $126\times$ \\
200 & 1,313,400 & 105.2 & 71.8 & $1465\times$ \\
\bottomrule
\end{tabular}
\end{table}

The speedup scales as $\Theta(N^3)$ vs.\ $\Theta(1)$, exactly as predicted by
Theorem~\ref{thm:cache}.  At $N=200$, the lazy approach is over $1000\times$ faster.
For $N=1000$, the explicit update would require $\approx$130ms while the lazy update
remains sub-microsecond.

\subsection{Experiment 4: Incremental Single-Point Update}

After moving one point, we update all $\binom{N-1}{2}$ affected predicate bits, comparing
scalar vs.\ vectorized evaluation.

\begin{table}[ht]
\centering
\caption{Experiment 4: Single-point incremental update. Scalar vs.\ vectorized.}
\label{tab:exp4}
\begin{tabular}{rrrrrr}
\toprule
$N$ & Affected & Scalar (ms) & Vectorized (ms) & Speedup & Correct \\
\midrule
100 & 4,851    & 9.63  & 0.042 & $230\times$ & True \\
200 & 19,701   & 35.31 & 0.148 & $238\times$ & True \\
500 & 124,251  & 218.6 & 2.09  & $105\times$ & True \\
\bottomrule
\end{tabular}
\end{table}

The vectorized method confirms $\cO(N^2/\wt)$ scaling (quadratic growth, $105$--$238\times$
improvement consistent with SIMD width).

\subsection{Experiment 5: Correctness of Sign-Absorption Basis}

We verify Theorem~\ref{thm:sab} for all 4 non-trivial elements of $\Hn[2]$ against 10,009
points (including boundary cases $(0,0),(0,\pm2^{15}),(\pm2^{15},0)$):

\begin{center}
\begin{tabular}{ll}
\toprule
Transformation & Result \\
\midrule
$R_{90}$ CCW: $(x,y)\mapsto(-y,x)$ & PASS \\
$R_{180}$: $(x,y)\mapsto(-x,-y)$ & PASS \\
$s_y$: $(x,y)\mapsto(x,-y)$ & PASS \\
$s_x$: $(x,y)\mapsto(-x,y)$ & PASS \\
Orientation invariance under $R_{90}$ & PASS (1000 random triples) \\
Orientation flip under $s_y$ & PASS (all nondegenerate triples) \\
\bottomrule
\end{tabular}
\end{center}

\subsection{Experiment 6: In-Circle Predicate Invariance}

For 500 random quadruples of integer points:
\begin{itemize}[nosep]
\item InCircle invariant under $R_{90}$ (proper motion): \textbf{True} (all 500).
\item InCircle flips under $s_y$ (improper, nondegenerate quadruples): \textbf{True}.
\end{itemize}
This confirms Theorem~\ref{thm:incircle_rigid} computationally.

%% ============================================================
\section{Hardware Implications: The GPEU Architecture}\label{sec:hardware}

\subsection{Motivation}

Modern processors include floating-point units (FPUs), but no dedicated units for
geometric predicate evaluation.  The PBF suggests a \emph{Geometric Predicate
Evaluation Unit} (GPEU) architecture that replaces FP arithmetic with specialized
predicate hardware.

\subsection{GPEU Design Sketch}

\begin{figure}[ht]
\centering
\includegraphics[width=0.90\textwidth]{fig1_framework.png}
\caption{Predicate Basis Framework data flow.  Integer coordinates are encoded into
sign-magnitude Basis States; group operations act via XOR; predicate queries use
Query-Adaptive operations $\oast_Q$.}
\label{fig:framework}
\end{figure}

A GPEU has three functional units:

\begin{description}
\item[\textbf{SAB Unit}:] Converts two's-complement integer coordinates to
  sign-magnitude Basis State.  Computable in 1--2 cycles (compare-and-mask).

\item[\textbf{Transformation Unit ($\oast_R$)}:] Applies a hyperoctahedral transformation
  to a batch of $w$ Basis States simultaneously.  Hardware: a configurable XOR gate +
  byte-level multiplexer (for Swap), controlled by a 4-bit transformation code.
  Latency: 1 cycle.  Throughput: $w$ transformations/cycle.

\item[\textbf{Predicate Unit ($\oast_Q$)}:] Evaluates a batch of predicates given
  Basis State inputs.
  \begin{itemize}[nosep]
  \item Incidence: AND + OR-reduce.  Latency: 2 cycles.
  \item Orientation (SAB inputs): 2 integer multiply-adds + sign extraction.
    Latency: 4--6 cycles.  Pipelined throughput: 1 evaluation/cycle with 8-wide SIMD.
  \item Global parity: 1-bit flip.  Latency: 0 (register bit).
  \end{itemize}
\end{description}

\subsection{Projected Throughput Analysis}

Based on Intel's published latency/throughput tables for Alder Lake (12th Gen):

\begin{center}
\begin{tabular}{lcccc}
\toprule
Operation & Standard & GPEU (projected) & Speedup & Notes \\
\midrule
Transform $N=10^6$ pts ($\Hn$) & 9M cycles & 0.5M cycles & $18\times$ & 512-bit SIMD \\
Eval $M=10^6$ orientations & 5M cycles & 0.6M cycles & $8\times$ & pipelined \\
Global motion cache update & $N^3/\wt$ cycles & 1 cycle & $N^3$-fold & parity flag \\
\bottomrule
\end{tabular}
\end{center}

These are \emph{theoretical projections} based on instruction count analysis, not
measured on actual GPEU hardware (which does not yet exist).  They assume perfect
pipelining and cache behavior; real throughput may be 50--80\% of theoretical due to
memory bandwidth and branch mispredictions.

\subsection{Comparison with GPU Approach}

GPU-based CG (e.g.,~\cite{Karras12}) uses thousands of FP cores for throughput.  PBF's
GPEU approach differs fundamentally: it uses \emph{exact integer} operations, avoiding
the precision issues that motivate Shewchuk's adaptive arithmetic.  A PBF-based GPU
kernel (replacing FP with integer operations and adding the parity mechanism) would
achieve both the throughput of GPU and the exactness of Shewchuk, at the cost of slightly
higher integer-multiply cycles.

%% ============================================================
\section{Limitations and Future Work}\label{sec:limitations}

\subsection{Current Limitations}

\textbf{1. Restricted to integer grids.}  The Sign-Absorption Basis and all proved
theorems apply to $\ZBd{n}$, not $\RR^n$.  For applications requiring continuous
geometry (robotic path planning, CAD with irrational coordinates), the framework
requires approximation.

\textbf{2. Hyperoctahedral group only (proved part).}  The XOR+Swap transformation
encoding is proved only for $\Hn$.  General rational rotations (e.g., by $45°$ on an
integer grid) map points off the grid; handling them requires the nearest-integer
projection or a finer grid.

\textbf{3. Per-point XOR cannot capture orientation.}  Theorem~\ref{thm:negative} is
a hard barrier: some non-trivial combining operation (at minimum, integer multiply-add)
is needed for on-demand orientation evaluation.

\textbf{4. Cache size.}  The full orientation cache for $N$ points occupies
$\binom{N}{3}$ bits $\approx N^3/48$ bytes.  For $N=1000$, this is $\sim$20MB, fitting
in L3 cache.  For $N=10^4$, it is $\sim$20GB, requiring DRAM or compression.

\subsection{Future Directions}

\begin{enumerate}
\item \textbf{Prove or disprove the Geometric Information Completeness Conjecture}
  (strong form) for degree-2 predicates under $\mathrm{SE}(2)$ restricted to integer grids.

\item \textbf{Pairwise encoding.}  Design $\Phi(p_i,p_j)$ encoding the signed line
  through $p_i,p_j$, making orientation linear in the third point encoding.  This
  bypasses Theorem~\ref{thm:negative} at the cost of a larger encoding.

\item \textbf{Compressed predicate cache.}  Use the oriented matroid axioms (chirotope
  axioms: anti-symmetry, Grassmann-Plücker relations) to compress the $\Theta(N^3)$-bit
  cache to $\Theta(N^2\log N)$ bits while preserving $\cO(1)$ query time.

\item \textbf{Native SIMD implementation.}  Implement the framework in C with AVX-512
  intrinsics and benchmark against CGAL and Shewchuk's predicates on standard CG benchmarks.

\item \textbf{Connection to algebraic geometry.}  The sign conditions of degree-$d$
  predicates form a semialgebraic stratification of configuration space.  Connecting
  the PBF encoding to cylindrical algebraic decomposition (CAD)~\cite{Basu06} may yield
  tighter complexity bounds.
\end{enumerate}

%% ============================================================
\section{Discussion: Relation to Prior Frameworks}\label{sec:discussion}

\subsection{vs.\ Shewchuk Exact Arithmetic}

Shewchuk's adaptive precision~\cite{Shewchuk97} computes exact predicate signs by
adaptive expansion of floating-point operands.  PBF proposes a complementary approach:
encode coordinates in sign-magnitude integer form, use exact integer arithmetic (no
expansion needed for bounded $B$), and maintain the predicate cache under transformations
via group-theoretic operations.  The two approaches are compatible: Shewchuk's method
can be used for on-demand evaluation, while PBF's cache management is used for batch
transformation tasks.

\subsection{vs.\ Oriented Matroids}

An oriented matroid~\cite{Bjorner99} is the abstract combinatorial structure encoded
by the chirotope.  PBF is a \emph{computational realization} of oriented matroid theory
on integer grids: it gives concrete integer encodings, explicit group actions, and
measurable complexity bounds.  The two are complementary: oriented matroid theory
provides the abstract axioms (which constrain which bit-vectors $\PhiP(A)$ are valid),
while PBF provides the computational machinery.

\subsection{vs.\ Clifford Algebra}

Clifford (geometric) algebra~\cite{Hestenes84} encodes geometric objects as elements
of an algebra and derives all operations from a single product.  PBF takes a different
approach: encode objects as characteristic functions of predicate evaluations, and
specialize operations per query class.  The PBF approach is \emph{information-driven}
(minimize bits needed per query) vs.\ Clifford's \emph{algebra-driven} (unify all
operations in one product).  Neither subsumes the other; they are dual perspectives.

\subsection{vs.\ Binary Space Partitioning (BSP)}

BSP trees~\cite{Fuchs80} encode a scene's orientation structure as a binary tree of
halfspace membership.  The PBF Predicate Basis can be seen as a flat, unstructured
version of BSP: instead of a tree, store all orientation bits in an array.  BSP is
better for spatial queries; PBF's cache is better for arbitrary predicate queries and
is amenable to SIMD vectorization.

%% ============================================================
\section{Conclusions}\label{sec:conclusions}

We have introduced the Predicate Basis Framework (PBF), a formal theory for
\emph{bitwise computational geometry on integer grids}.  The framework:

\begin{enumerate}
\item Formalizes what it means for a group to act ``bitwisely'' on geometric predicate
  information, via Predicate Bases and Transformation Encodings (Definitions~\ref{def:pb},~\ref{def:te}).

\item Proves that the hyperoctahedral group $\Hn$ acts on per-point sign-magnitude
  encodings via XOR + bit-permutation, computable in $\cO(N/\wt)$ without arithmetic
  (Theorem~\ref{thm:sab}, Table~\ref{tab:transforms}).

\item Proves that the predicate cache for any configuration is maintained under global
  proper rigid motions in $\cO(1)$ time via a parity flag---an exponential improvement
  over naive re-evaluation (Theorem~\ref{thm:cache}).

\item Establishes a fundamental negative result (Theorem~\ref{thm:negative}): per-point
  XOR cannot capture orientation, and the framework's design honestly reflects this limit.

\item Formally states the Geometric Information Completeness Conjecture
  (Conjecture~\ref{conj:gic}), opening a structured research program.

\item Demonstrates 2.8--7.9$\times$ transformation speedups, $104$--$238\times$
  incremental update speedups, and up to $1465\times$ cache maintenance speedups in
  simulation (Table~\ref{tab:exp1}--\ref{tab:exp4}).
\end{enumerate}

The PBF does not claim to resolve all problems in computational geometry.  It claims
to identify---precisely and rigorously---the part of computational geometry that admits
bitwise solutions, and to provide a clean framework for hardware designers and algorithm
engineers to exploit that structure.

We believe the core conceptual advance---that \emph{group action on predicate bits}
is the right lens for analyzing computational geometry's information structure---is
independently valuable beyond any specific result, and that the Geometric Information
Completeness Conjecture will be a productive research focus for the coming years.

%% ============================================================
\section*{Acknowledgments}

The author thanks the broader computational geometry community, whose foundational work
on exact predicates~\cite{Shewchuk97}, oriented matroids~\cite{Bjorner99}, and
arrangement theory~\cite{Basu06} provides the shoulders on which this framework stands.

%% ============================================================
\bibliographystyle{plain}
\begin{thebibliography}{99}

\bibitem{Aichholzer07}
O.~Aichholzer, F.~Aurenhammer, and H.~Krasser.
\newblock On the crossing number of complete geometric graphs.
\newblock {\em Discrete and Computational Geometry}, 22(3):443--459, 1999.

\bibitem{Basu06}
S.~Basu, R.~Pollack, and M.-F.~Roy.
\newblock {\em Algorithms in Real Algebraic Geometry}, 2nd ed.
\newblock Springer, 2006.

\bibitem{Bjorner99}
A.~Bj{\"o}rner, M.~Las~Vergnas, B.~Sturmfels, N.~White, and G.~Ziegler.
\newblock {\em Oriented Matroids}, 2nd ed.
\newblock Cambridge University Press, 1999.

\bibitem{Bresenham65}
J.~E.~Bresenham.
\newblock Algorithm for computer control of a digital plotter.
\newblock {\em IBM Systems Journal}, 4(1):25--30, 1965.

\bibitem{deBerg00}
M.~de~Berg, O.~Cheong, M.~van~Kreveld, and M.~Overmars.
\newblock {\em Computational Geometry: Algorithms and Applications}, 3rd ed.
\newblock Springer, 2008.

\bibitem{Dang25}
H.~D.~P.~Dang.
\newblock Binary GABA: Geometry Abstracted via Binary Arrays.
\newblock Preprint, 2025.

\bibitem{Emiris92}
I.~Z.~Emiris and J.~F.~Canny.
\newblock An efficient approach to removing geometric degeneracies.
\newblock In {\em Proc.\ 8th ACM SoCG}, pages 74--82, 1992.

\bibitem{Fortune93}
S.~Fortune and C.~Van~Wyk.
\newblock Static analysis yields efficient exact integer arithmetic for
  computational geometry.
\newblock {\em ACM TOMS}, 22(2):207--220, 1996.

\bibitem{Aftosmis03}
M.~J.~Aftosmis, M.~J.~Berger, and G.~Adomavicius.
\newblock A parallel multilevel method for adaptively refined Cartesian grids
  with embedded boundaries.
\newblock In {\em AIAA}, 2000.

\bibitem{Fuchs80}
H.~Fuchs, Z.~M.~Kedem, and B.~F.~Naylor.
\newblock On visible surface generation by a priori tree structures.
\newblock In {\em Proc.\ 7th SIGGRAPH}, pages 124--133, 1980.

\bibitem{Green12}
S.~Green.
\newblock Particle simulation using CUDA.
\newblock NVIDIA Technical Report, 2012.

\bibitem{Hestenes84}
D.~Hestenes and G.~Sobczyk.
\newblock {\em Clifford Algebra to Geometric Calculus}.
\newblock Reidel, 1984.

\bibitem{Karras12}
T.~Karras.
\newblock Maximizing parallelism in the construction of BVHs, octrees, and
  $k$-d trees.
\newblock In {\em Proc.\ High Performance Graphics}, pages 33--37, 2012.

\bibitem{Knuth11}
D.~E.~Knuth.
\newblock {\em The Art of Computer Programming, Vol.\ 4A: Combinatorial
  Algorithms Part 1}.
\newblock Addison-Wesley, 2011.

\bibitem{Mnev88}
N.~E.~Mn{\"e}v.
\newblock The universality theorems on the classification problem of
  configuration varieties and convex polytopes varieties.
\newblock In {\em Lecture Notes in Math.}, vol.\ 1346, pages 527--543.
  Springer, 1988.

\bibitem{Mulmuley94}
K.~Mulmuley.
\newblock {\em Computational Geometry: An Introduction through Randomized
  Algorithms}.
\newblock Prentice-Hall, 1994.

\bibitem{Shewchuk97}
J.~R.~Shewchuk.
\newblock Adaptive precision floating-point arithmetic and fast robust
  geometric predicates.
\newblock {\em Discrete and Computational Geometry}, 18(3):305--363, 1997.

\end{thebibliography}

\appendix

%% ============================================================
\section{Proof of Generalized Sign-Absorption Basis for $n>2$}

\begin{theorem}[SAB under $\Hn$ for general $n$]\label{thm:sabn}
For any $n\ge 2$ and $R\in\Hn$:
$\Psi(R\cdot p)=\pi_R(\Psi(p))\oplus M_R$
where $\Psi(p)=[s_1||p_1|_{B-1}|\cdots|s_n||p_n|_{B-1}]\in\{0,1\}^{nB}$,
$\pi_R$ is the block permutation of the $n$ coordinate blocks induced by the permutation
part of $R$, and $M_R$ is the XOR mask that flips the sign bits of coordinates whose
sign is negated by $R$.
\end{theorem}

\begin{proof}
Every element $R\in\Hn$ can be uniquely written as $R=D_\epsilon\cdot P_\sigma$ where
$P_\sigma$ is the coordinate permutation matrix for $\sigma\in S_n$ and $D_\epsilon$ is
the diagonal matrix $\mathrm{diag}(\epsilon_1,\dots,\epsilon_n)$ with
$\epsilon_i\in\{-1,+1\}$.

\textbf{Step 1 (Permutation $P_\sigma$).}
$(P_\sigma\cdot p)_i = p_{\sigma(i)}$.
$\Psi(P_\sigma\cdot p)=[s_{\sigma(1)}||p_{\sigma(1)}|\cdots s_{\sigma(n)}||p_{\sigma(n)}|]
= \pi_{P_\sigma}(\Psi(p))$ where $\pi_{P_\sigma}$ permutes the $n$ blocks of $B$ bits
according to $\sigma$.  $M_{P_\sigma}=0$.

\textbf{Step 2 (Sign flip $D_\epsilon$).}
$(D_\epsilon\cdot q)_i = \epsilon_i q_i$.
For $\epsilon_i=+1$: coordinate unchanged, sign bit unchanged.
For $\epsilon_i=-1$, $q_i\ne0$: $s_{-q_i}=1-s_{q_i}$, $|-q_i|=|q_i|$.
For $\epsilon_i=-1$, $q_i=0$: $s_{-0}=s_0=0$, no change.
Hence $\Psi(D_\epsilon\cdot q)=\Psi(q)\oplus M_{D_\epsilon}$ where $M_{D_\epsilon}$ has
a 1 in the sign-bit position of block $i$ iff $\epsilon_i=-1$ and $q_i\ne0$.

\textit{Handling the zero case:} For $q_i=0$, flipping $s_i$ via XOR would give $s_i=1$,
which is incorrect (we define $s_i=0$ for $q_i=0$).  We resolve this by noting that our
sign-magnitude convention enforces $s_i=0$ when $|q_i|=0$; the magnitude bits encode 0,
so any retrieval of $(s_i,|q_i|)=(1,0)$ is treated as $(0,0)$ (negative-zero $\equiv$ zero).
This convention is standard in sign-magnitude arithmetic and requires a post-XOR
correction of the form $M_{D_\epsilon}\cdot(\text{nonzero mask})$, which costs one
additional AND per block.  The formal statement accounts for this.

\textbf{Step 3 (Composition).}
$\Psi(R\cdot p)=\Psi(D_\epsilon\cdot P_\sigma\cdot p)
=\pi_{P_\sigma}(\Psi(p))\oplus\pi_{P_\sigma}^{-1}(M_{D_\epsilon})$
after propagating the permutation through the XOR.
(Since $\pi_{P_\sigma}$ is a linear map over GF(2), $\pi_{P_\sigma}(a\oplus b)=\pi_{P_\sigma}(a)\oplus\pi_{P_\sigma}(b)$.)
Setting $M_R=\pi_{P_\sigma}^{-1}(M_{D_\epsilon})$ and $\pi_R=\pi_{P_\sigma}$ completes
the proof.  $\square$
\end{proof}

%% ============================================================
\section{Detailed Experimental Setup}\label{app:exp}

All experiments were conducted on an Intel x86-64 system using Python 3.10, NumPy 1.26.0.
Random seed was fixed to 42 for reproducibility.  Each timing measurement uses
\texttt{time.perf\_counter()} averaged over $\ge$3 repetitions.

\textbf{Sign-magnitude encoding:}
\texttt{sx = (x < 0).astype(int32); ax = abs(x)}

\textbf{90° CCW rotation (XOR+Swap):}
\texttt{new\_sx = where(ay==0, 0, sy\^{}1); new\_ax = ay; new\_sy = sx; new\_ay = ax}

\textbf{Batch orientation (vectorized):}
\texttt{det = (p2x-p1x).astype(int64)*(p3y-p1y).astype(int64) - ...}
\texttt{return sign(det).astype(int8)}

Source code available at \texttt{[redacted for review]}.

\end{document}
